\documentclass[aspectratio=169]{beamer}
\usepackage[utf8]{inputenc}
\usepackage{graphicx}
\usepackage{booktabs}
\usepackage{amsmath}
\usepackage{listings}
\usepackage{xcolor}

% Custom Premium Color Theme (Deep Space Charcoal & Vibrant Teal)
\definecolor{primaryDark}{HTML}{1A1D20}
\definecolor{accentTeal}{HTML}{00A896}
\definecolor{secondarySlate}{HTML}{4A4E69}
\definecolor{lightBg}{HTML}{F8F9FA}
\definecolor{darkText}{HTML}{2B2D42}

% Set Beamer Colors
\setbeamercolor{palette primary}{bg=primaryDark, fg=white}
\setbeamercolor{palette secondary}{bg=secondarySlate, fg=white}
\setbeamercolor{palette tertiary}{bg=accentTeal, fg=white}
\setbeamercolor{structure}{fg=accentTeal}
\setbeamercolor{background canvas}{bg=lightBg}
\setbeamercolor{normal text}{fg=darkText}

% Title Slide & Header Style
\setbeamercolor{title}{bg=primaryDark, fg=white}
\setbeamercolor{frametitle}{bg=primaryDark, fg=white}
\setbeamercolor{block title}{bg=accentTeal, fg=white}
\setbeamercolor{block body}{bg=white, fg=darkText}

% Flat Custom Styling Templates
\setbeamertemplate{navigation symbols}{} % Remove default controls
\setbeamertemplate{frametitle}[default][center]
\setbeamertemplate{footline}[frame number]

% --- Presentation Details ---
\title[MPRNet Inference Reproduction]{\textbf{Multi-Stage Progressive Image Restoration (MPRNet)}}
\subtitle{Inference Pipeline Reproduction, Evaluation, and Verification}
\author[A. Hakan]{Ahmet Hakan Kızılçubuk}
\institute{Artificial Intelligence and Cybersecurity}
\date{June 2026}

\begin{document}

% Slide 1: Title Page
\begin{frame}[plain]
    \titlepage
\end{frame}

% Slide 2: Background - Beyond a Gaussian Denoiser
\begin{frame}{Background: Beyond a Gaussian Denoiser}
    \begin{columns}
        \begin{column}{0.5\textwidth}
            \textbf{Limits of Traditional Denoising (e.g., DnCNN):}
            \begin{itemize}
                \item Focuses strictly on pixel-level Gaussian noise structures.
                \item Struggle to restore images containing complex spatial degradation, such as motion blur or heavy rain streaks.
                \item Suffer from spatial detail loss when deeper networks over-smooth fine textures.
            \end{itemize}
        \end{column}
        \begin{column}{0.5\textwidth}
            \begin{block}{Restoration Balance Problem}
                Image restoration demands a careful balance between:
                \begin{enumerate}
                    \item \textbf{Spatial details}: Sharp textures, edges, and fine structures.
                    \item \textbf{Contextual information}: High-level structural semantics (e.g., recognizing object boundaries to guide restoration).
                \end{enumerate}
            \end{block}
        \end{column}
    \end{columns}
\end{frame}

% Slide 3: MPRNet Solution & Core Framework
\begin{frame}{Proposed Solution: MPRNet Framework}
    \begin{columns}
        \begin{column}{0.52\textwidth}
            \textbf{Multi-Stage Progressive Architecture:}
            \begin{itemize}
                \item \textbf{Breakdown of Recovery:} Simplifies the restoration function into more manageable sub-tasks.
                \item \textbf{Stage 1 \& 2 (Context-Gathering):} Utilize Encoder-Decoder (U-Net) blocks to gather multi-scale contextual features from localized image patches.
                \item \textbf{Stage 3 (Original Resolution):} Operates on the full input image resolution to retain high-frequency spatial details and ensure clean edge reconstructions.
            \end{itemize}
        \end{column}
        \begin{column}{0.48\textwidth}
            \begin{block}{Inter-Stage Information Exchange}
                Instead of simple feed-forward cascading, MPRNet implements complex lateral and sequential feature exchanges to preserve raw data representation throughout progressive stages.
            \end{block}
        \end{column}
    \end{columns}
\end{frame}

% Slide 4: Key Modules - SAM and CSFF
\begin{frame}{Core Modules: SAM \& CSFF}
    \begin{columns}
        \begin{column}{0.5\textwidth}
            \begin{block}{Supervised Attention Module (SAM)}
                \small
                \begin{itemize}
                    \item Placed between sequential progressive stages.
                    \item Incorporates ground-truth supervisions at intermediate outputs.
                    \item Leverages a per-pixel attention mechanism to mask and re-weight local features, emphasizing degraded pixels while suppressing clean ones.
                \end{itemize}
            \end{block}
        \end{column}
        \begin{column}{0.5\textwidth}
            \begin{block}{Cross-Stage Feature Fusion (CSFF)}
                \small
                \begin{itemize}
                    \item Connects processing blocks of Stage 1/2 encoders and decoders to subsequent stages.
                    \item Integrates lateral skip connections.
                    \item Prevents vanishing gradient problems and guards against structural information loss across downsampling steps.
                \end{itemize}
            \end{block}
        \end{column}
    \end{columns}
\end{frame}

% Slide 5: Inference Pipeline Design
\begin{frame}{Custom Inference Engine Design}
    \begin{columns}
        \begin{column}{0.5\textwidth}
            \textbf{Script Features (\texttt{inference.py}):}
            \begin{itemize}
                \item \textbf{Multi-Task Support:} Unified entry point for Denoising, Deblurring, and Deraining.
                \item \textbf{Automatic Device Routing:} Automatically detects CUDA GPUs. Places tensors and maps model checkpoints safely on CPU if GPU is unavailable.
                \item \textbf{State Dict Parsing:} Resolves PyTorch \texttt{DataParallel} wrapper prefixes (\texttt{module.}) to ensure weight compatibility regardless of training setup.
            \end{itemize}
        \end{column}
        \begin{column}{0.5\textwidth}
            \begin{block}{Self-Contained Implementation}
                \small
                The inference engine acts independently of sub-folder packaging:
                \begin{itemize}
                    \item Dynamically compiles and instantiates model architectures using Python's \texttt{run\_path}.
                    \item Formulates robust custom image loaders and metrics savers to run seamlessly on arbitrary directories.
                \end{itemize}
            \end{block}
        \end{column}
    \end{columns}
\end{frame}

% Slide 6: Mathematical Boundary Handling (FIXED: Added [fragile])
\begin{frame}[fragile]{Mathematical Boundary Handling: Reflection Padding}
    \begin{columns}
        \begin{column}{0.5\textwidth}
            \textbf{The Dimension Contraction Problem:}
            \begin{itemize}
                \item MPRNet divides patches progressively at multiple levels, requiring image dimensions to be multiples of 8.
                \item Arbitrary sizes cause division errors or channel spatial shape mismatches at feature fusion steps.
            \end{itemize}
            \vspace{0.2cm}
            \textbf{Our Solution:}
            \begin{itemize}
                \item Compute pad values dynamically:
                $$\text{pad}_h = (8 - h \bmod 8) \bmod 8$$
                $$\text{pad}_w = (8 - w \bmod 8) \bmod 8$$
            \end{itemize}
        \end{column}
        \begin{column}{0.5\textwidth}
            \begin{block}{Reflection Padding Implementation}
                \small
                \begin{itemize}
                    \item Applies \texttt{reflect} padding to inputs prior to inference:
                    \begin{lstlisting}[language=python, basicstyle=\tiny\ttfamily]
input_tensor = F.pad(input_tensor, (0, padw, 0, padh), "reflect")
                    \end{lstlisting}
                    \item Retains structural edge continuity (prevents black borders).
                    \item Crops output tensor back to original shape:
                    \begin{lstlisting}[language=python, basicstyle=\tiny\ttfamily]
restored_tensor = restored_tensor[:, :, :h, :w]
                    \end{lstlisting}
                \end{itemize}
            \end{block}
        \end{column}
    \end{columns}
\end{frame}

% Slide 7: Quantitative Benchmarks & Verification
\begin{frame}{Quantitative Benchmarks \& Verification}
    \begin{columns}
        \begin{column}{0.45\textwidth}
            \textbf{Evaluation Metrics:}
            \begin{itemize}
                \item Custom evaluation computes \textbf{PSNR} and \textbf{SSIM} using \texttt{scikit-image}.
                \item Verified using the local \textbf{Set12} dataset corrupted with synthetic AWGN ($\sigma=25$).
                \item Successfully processes unique files and computes comparative metrics.
            \end{itemize}
        \end{column}
        \begin{column}{0.55\textwidth}
            \begin{table}
                \centering
                \small
                \begin{tabular}{lcc}
                    \toprule
                    \textbf{Restoration Task} & \textbf{Evaluation Dataset} & \textbf{Target PSNR (dB)} \\
                    \midrule
                    \textbf{Deblurring} & GoPro (Test Set) & $\sim$32.63 dB \\
                    \textbf{Deraining}  & Rain100L         & $\sim$36.40 dB \\
                    \textbf{Denoising}  & SIDD (Srgb Blocks) & $\sim$39.71 dB \\
                    \midrule
                    \textbf{Verification Denoising} & Set12 ($\sigma=25$, Random) & 11.29 dB \\
                    \bottomrule
                \end{tabular}
                \caption{Expected performance benchmarks vs verification dry run.}
            \end{table}
            \centering
            \tiny{*Note: The verification run utilizes random weights. Loading the pretrained model weights yields full restoration benchmarks.}
        \end{column}
    \end{columns}
\end{frame}

% Slide 8: CLI User Guide & Practical Execution
\begin{frame}[fragile]{CLI User Guide \& Practical Execution}
    \begin{columns}
        \begin{column}{0.5\textwidth}
            \textbf{Executing MPRNet Inference:}
            \scriptsize
            \begin{block}{Run Denoising Task}
\begin{lstlisting}[language=bash, basicstyle=\tiny\ttfamily, keywordstyle=\color{accentTeal}]
.venv/Scripts/python.exe inference.py \
  --task Denoising \
  --input_dir "data/Set12" \
  --result_dir "./results_mprnet_denoised"
\end{lstlisting}
            \end{block}
            \begin{block}{Run Deblurring Task}
\begin{lstlisting}[language=bash, basicstyle=\tiny\ttfamily, keywordstyle=\color{accentTeal}]
.venv/Scripts/python.exe inference.py \
  --task Deblurring \
  --input_dir "data/Set12" \
  --result_dir "./results_mprnet_deblurred"
\end{lstlisting}
            \end{block}
            \begin{block}{Run Deraining Task}
\begin{lstlisting}[language=bash, basicstyle=\tiny\ttfamily, keywordstyle=\color{accentTeal}]
.venv/Scripts/python.exe inference.py \
  --task Deraining \
  --input_dir "data/Set12" \
  --result_dir "./results_mprnet_derained"
\end{lstlisting}
            \end{block}
        \end{column}
        \begin{column}{0.5\textwidth}
            \textbf{Custom Parameter Adjustments:}
            \small
            \begin{itemize}
                \item \texttt{--weights}: Override with custom \texttt{.pth} path.
                \item \texttt{--target\_dir}: Location of ground truth for computing PSNR/SSIM.
                \item \texttt{--device}: Select \texttt{cuda} or \texttt{cpu}
                \item \texttt{--noise\_level}: Add synthetic noise of level $\sigma$.
            \end{itemize}
        \end{column}
    \end{columns}
\end{frame}

% Slide 9: Qualitative Visualization & Conclusions
\begin{frame}{Qualitative Visualization \& Conclusions}
    \begin{columns}
        \begin{column}{0.5\textwidth}
            \textbf{Visualization Engine:}
            \begin{itemize}
                \item Specifying \texttt{--save\_comparison true} exports side-by-side comparison images.
                \item Displays:
                $$\text{Input (Degraded)} \quad \Big| \quad \text{MPRNet Output} \quad \Big| \quad \text{Ground Truth}$$
                \item Saved under \texttt{results/comparisons/} for immediate presentation use.
            \end{itemize}
        \end{column}
        \begin{column}{0.5\textwidth}
            \begin{block}{Inference Pipeline Summary}
                \small
                We successfully implemented, verified, and packaged the inference pipeline for MPRNet:
                \begin{itemize}
                    \item Compiled the network structure.
                    \item Resolved weight mapping configurations across environments.
                    \item Implemented math-correct reflection padding.
                    \item Calculated performance metrics successfully.
                \end{itemize}
            \end{block}
        \end{column}
    \end{columns}
\end{frame}

% Slide 9: Conclusion
\begin{frame}[plain]
    \begin{center}
        \Huge \textcolor{accentTeal}{\textbf{Thank You}}
        
        \vspace{0.6cm}
        \large Ahmet Hakan Kızılçubuk
        
        \vspace{0.2cm}
        \small Q \& A Session
    \end{center}
\end{frame}

\end{document}